
\medskip
\noindent
\textbf{Triangle coherence.}
The three functors satisfy
$\gamma = \alpha^{-1} \circ \beta^{-1}$. Each functor
is an equivalence with explicit inverse as constructed
above, and the triangle~\eqref{eq:equiv-triangle}
commutes: $\alpha \circ \gamma \circ \beta = \id$.
\end{proof}

\begin{remark}[Independent proof via the cohomological Hall algebra]
\label{rem:independent-proof-coha}
The equivalence $\textup{(I)} \leftrightarrow \textup{(II)}
\leftrightarrow \textup{(III)}$ admits a second, independent
proof for algebras arising as critical CoHAs of quivers with
potential, using the framework of
Kontsevich--Soibelman~\cite{KS11},
Schiffmann--Vasserot~\cite{SV13},
Yang--Zhao~\cite{YZ18a}, and
Jindal--Kaubrys--Latyntsev~\cite{JKL26}. The logic is
orthogonal to the bar-cobar adjunction above: rather than
constructing the triangle abstractly from the ordered bar
complex, one exhibits all three structures simultaneously on a
single geometric object and verifies that they determine each
other.

\smallskip
\noindent\textbf{Step~1: simultaneous existence of (I) and (III)
on the CoHA.}
Let $(Q,W)$ be a quiver with potential. The critical CoHA
\[
  \cH_{Q,W}
  = \bigoplus_{d \in \ZZ_{\geq 0}^{Q_0}}
  H^{G_d}_{\mathrm{BM},*}
  (\cM_d^{\mathrm{nil}},\, \varphi_W)
\]
carries two structures of independent geometric origin:
\begin{enumerate}[label=\textup{(\alph*)}]
\item \emph{Vertex coproduct} $\Delta_{\mathrm{CoHA}}$, from
  restriction to Levi subalgebras
  $G_{d_1} \times G_{d_2} \hookrightarrow G_d$; and
\item \emph{$R$-matrix} $R(z)$, from the
  Maulik--Okounkov stable envelope
  construction~\cite{MO19} on the Nakajima variety
  associated to~$Q$.
\end{enumerate}
Jindal--Kaubrys--Latyntsev~\cite{JKL26} prove that these two
structures satisfy the vertex bialgebra axiom (their
Theorems~A and~C):
\begin{equation}\label{eq:jkl-vertex-bialgebra}
  \Delta_{\mathrm{CoHA},z}(Y(a,w)\,b)
  = Y^{(2)}(\Delta_{\mathrm{CoHA},z}(a),w)\,
  \Delta_{\mathrm{CoHA},z}(b),
\end{equation}
which is the OPE compatibility
axiom~\eqref{eq:ope-compat}. The $R$-matrix satisfies the
QYBE and unitarity by the stable envelope formalism
\cite{MO19}; the quasi-coassociativity of
$\Delta_{\mathrm{CoHA}}$ follows from the
associativity of the Hall product on the Drinfeld
double~\cite{YZ18a}. This establishes the equivalence
$\textup{(I)} \leftrightarrow \textup{(III)}$ by
construction: the $R$-matrix and the coproduct coexist on
$\cH_{Q,W}$ and satisfy the quasi-triangularity relation
$R(z)\, \Delta^{\mathrm{ch}} = \Delta^{\mathrm{ch},\op}\,
R(z)$.

\smallskip
\noindent\textbf{Step~2: extraction of (II) from the CoHA
vertex bialgebra.}
The CoHA $\cH_{Q,W}$, equipped with its $R$-matrix, is an
$\Eone$-chiral algebra in the sense of
Definition~\ref{def:e1-chiral-rmatrix}. The functor $\alpha$
of the equivalence triangle~\eqref{eq:equiv-triangle}
extracts the chiral $\Ainf$ operations
$\{m_k^{\mathrm{ch}}\}$ by restricting the OPE to the
associahedra
$K_k \hookrightarrow
\overline{\FM}_k^{\mathrm{ord}}(\CC)$ as
in~\eqref{eq:mk-from-ope}. The Stasheff identities follow
from Stokes' theorem on $K_{k+1}$ as before; the
convergence of the integrals is guaranteed by
$S$-locality of the $R$-matrix. This gives
$\textup{(I)} \to \textup{(II)}$ on the CoHA.

The reverse direction $\textup{(II)} \to \textup{(I)}$ is the
degree-$2$ holonomy recovery of the $R$-matrix from
$m_2^{\mathrm{ch}}$, which depends only on the formal
structure of the Stasheff associahedra and not on the CoHA
origin.

\smallskip
\noindent\textbf{Step~3: the triangle closes.}
By the JKL vertex bialgebra theorem, the CoHA
simultaneously carries (I), (II), and (III). To verify that
these three structures \emph{determine} each other (not merely
coexist), one checks the uniqueness directions. On the Koszul
locus:
\begin{enumerate}[label=\textup{(\roman*)}]
\item \emph{(I) determines (III):} given an $R$-matrix
  satisfying the QYBE and $S$-locality, the quasi-triangular
  coproduct is unique up to gauge equivalence by the
  Drinfeld reconstruction~\cite{Drinfeld90}
  (spectral parameter version:
  Etingof--Kazhdan~\cite{EK96}). The JKL coproduct on the
  CoHA is the unique one compatible with the MO $R$-matrix.
\item \emph{(III) determines (I):} the Drinfeld
  formula~\eqref{eq:r-from-coprod-proof} recovers $S(z)$
  from $\Delta^{\mathrm{ch}}$, independently of the CoHA
  construction.
\item \emph{(I) $\leftrightarrow$ (II):} the extraction of
  $\Ainf$ operations from the $R$-matrix is the
  Stasheff--FM construction (direction $\alpha$), whose
  inverse recovers the $R$-matrix as holonomy (direction
  $\alpha^{-1}$). Neither direction uses the CoHA.
\end{enumerate}
The triangle coherence
$\alpha \circ \gamma \circ \beta = \id$ then follows from
the uniqueness of each vertex.

\smallskip
\noindent\textbf{Independence from the bar-cobar proof.}
The argument above differs from the proof of
Theorem~\ref{thm:chiral-qg-equiv} in two respects.
First, the equivalence $\textup{(I)} \leftrightarrow
\textup{(III)}$ is established by the geometric CoHA
construction (stable envelopes for the $R$-matrix,
Levi restriction for the coproduct, JKL for the vertex
bialgebra axiom), bypassing the ordered bar complex
$\Barord(\cA) = T^c(s^{-1}\bar{\cA})$ and the
cobar dualization~\eqref{eq:coprod-from-bar}. Second, the
uniqueness of the coproduct given the $R$-matrix is
deduced from the Drinfeld reconstruction theorem for
quasi-triangular vertex bialgebras, rather than from the
bar-cobar adjunction on the Koszul locus. The functor
$\alpha$ (extraction of $\Ainf$ operations from the
$R$-matrix via associahedron integrals) is shared between
